% Options for packages loaded elsewhere
\PassOptionsToPackage{unicode}{hyperref}
\PassOptionsToPackage{hyphens}{url}
\PassOptionsToPackage{dvipsnames,svgnames,x11names}{xcolor}
%
\documentclass[
  12pt]{article}

\usepackage{amsmath,amssymb}
\usepackage{iftex}
\ifPDFTeX
  \usepackage[T1]{fontenc}
  \usepackage[utf8]{inputenc}
  \usepackage{textcomp} % provide euro and other symbols
\else % if luatex or xetex
  \usepackage{unicode-math}
  \defaultfontfeatures{Scale=MatchLowercase}
  \defaultfontfeatures[\rmfamily]{Ligatures=TeX,Scale=1}
\fi
\usepackage{lmodern}
\ifPDFTeX\else  
    % xetex/luatex font selection
\fi
% Use upquote if available, for straight quotes in verbatim environments
\IfFileExists{upquote.sty}{\usepackage{upquote}}{}
\IfFileExists{microtype.sty}{% use microtype if available
  \usepackage[]{microtype}
  \UseMicrotypeSet[protrusion]{basicmath} % disable protrusion for tt fonts
}{}
\makeatletter
\@ifundefined{KOMAClassName}{% if non-KOMA class
  \IfFileExists{parskip.sty}{%
    \usepackage{parskip}
  }{% else
    \setlength{\parindent}{0pt}
    \setlength{\parskip}{6pt plus 2pt minus 1pt}}
}{% if KOMA class
  \KOMAoptions{parskip=half}}
\makeatother
\usepackage{xcolor}
\setlength{\emergencystretch}{3em} % prevent overfull lines
\setcounter{secnumdepth}{5}
% Make \paragraph and \subparagraph free-standing
\makeatletter
\ifx\paragraph\undefined\else
  \let\oldparagraph\paragraph
  \renewcommand{\paragraph}{
    \@ifstar
      \xxxParagraphStar
      \xxxParagraphNoStar
  }
  \newcommand{\xxxParagraphStar}[1]{\oldparagraph*{#1}\mbox{}}
  \newcommand{\xxxParagraphNoStar}[1]{\oldparagraph{#1}\mbox{}}
\fi
\ifx\subparagraph\undefined\else
  \let\oldsubparagraph\subparagraph
  \renewcommand{\subparagraph}{
    \@ifstar
      \xxxSubParagraphStar
      \xxxSubParagraphNoStar
  }
  \newcommand{\xxxSubParagraphStar}[1]{\oldsubparagraph*{#1}\mbox{}}
  \newcommand{\xxxSubParagraphNoStar}[1]{\oldsubparagraph{#1}\mbox{}}
\fi
\makeatother


\providecommand{\tightlist}{%
  \setlength{\itemsep}{0pt}\setlength{\parskip}{0pt}}\usepackage{longtable,booktabs,array}
\usepackage{calc} % for calculating minipage widths
% Correct order of tables after \paragraph or \subparagraph
\usepackage{etoolbox}
\makeatletter
\patchcmd\longtable{\par}{\if@noskipsec\mbox{}\fi\par}{}{}
\makeatother
% Allow footnotes in longtable head/foot
\IfFileExists{footnotehyper.sty}{\usepackage{footnotehyper}}{\usepackage{footnote}}
\makesavenoteenv{longtable}
\usepackage{graphicx}
\makeatletter
\def\maxwidth{\ifdim\Gin@nat@width>\linewidth\linewidth\else\Gin@nat@width\fi}
\def\maxheight{\ifdim\Gin@nat@height>\textheight\textheight\else\Gin@nat@height\fi}
\makeatother
% Scale images if necessary, so that they will not overflow the page
% margins by default, and it is still possible to overwrite the defaults
% using explicit options in \includegraphics[width, height, ...]{}
\setkeys{Gin}{width=\maxwidth,height=\maxheight,keepaspectratio}
% Set default figure placement to htbp
\makeatletter
\def\fps@figure{htbp}
\makeatother

\addtolength{\oddsidemargin}{-.5in}%
\addtolength{\evensidemargin}{-.1in}%
\addtolength{\textwidth}{1in}%
\addtolength{\textheight}{1.7in}%
\addtolength{\topmargin}{-1in}
\makeatletter
\@ifpackageloaded{caption}{}{\usepackage{caption}}
\AtBeginDocument{%
\ifdefined\contentsname
  \renewcommand*\contentsname{Table of contents}
\else
  \newcommand\contentsname{Table of contents}
\fi
\ifdefined\listfigurename
  \renewcommand*\listfigurename{List of Figures}
\else
  \newcommand\listfigurename{List of Figures}
\fi
\ifdefined\listtablename
  \renewcommand*\listtablename{List of Tables}
\else
  \newcommand\listtablename{List of Tables}
\fi
\ifdefined\figurename
  \renewcommand*\figurename{Figure}
\else
  \newcommand\figurename{Figure}
\fi
\ifdefined\tablename
  \renewcommand*\tablename{Table}
\else
  \newcommand\tablename{Table}
\fi
}
\@ifpackageloaded{float}{}{\usepackage{float}}
\floatstyle{ruled}
\@ifundefined{c@chapter}{\newfloat{codelisting}{h}{lop}}{\newfloat{codelisting}{h}{lop}[chapter]}
\floatname{codelisting}{Listing}
\newcommand*\listoflistings{\listof{codelisting}{List of Listings}}
\makeatother
\makeatletter
\makeatother
\makeatletter
\@ifpackageloaded{caption}{}{\usepackage{caption}}
\@ifpackageloaded{subcaption}{}{\usepackage{subcaption}}
\makeatother


\ifLuaTeX
  \usepackage{selnolig}  % disable illegal ligatures
\fi
\usepackage[]{natbib}
\bibliographystyle{agsm}
\usepackage{bookmark}

\IfFileExists{xurl.sty}{\usepackage{xurl}}{} % add URL line breaks if available
\urlstyle{same} % disable monospaced font for URLs
\hypersetup{
  pdftitle={Title},
  pdfauthor={Scott LaCombe},
  pdfkeywords={3 to 6 keywords, that do not appear in the title},
  colorlinks=true,
  linkcolor={blue},
  filecolor={Maroon},
  citecolor={Blue},
  urlcolor={Blue},
  pdfcreator={LaTeX via pandoc}}




\newcommand{\anon}{1}

%set the key \texttt{anon} to ``0'' to hide the authors and acknowledgements,
%  producing the required anonymized version. 
%Set the key \texttt{anon} to ``1'' to produce the manuscript with author details and
% acknowledgments. 



\begin{document}


\def\spacingset#1{\renewcommand{\baselinestretch}%
{#1}\small\normalsize} \spacingset{1}


%%%%%%%%%%%%%%%%%%%%%%%%%%%%%%%%%%%%%%%%%%%%%%%%%%%%%%%%%%%%%%%%%%%%%%%%%%%%%%


\if1\anon
{
  \title{\bf Teaching Reproducibility and Replicability: \\Principles and Examples}
  \author{
  Scott J Lacombe\thanks{This work stems from a workshop with grant support from The Alliance to Advance Liberal Arts Colleges.}\hspace{.2cm} ORCID:0000-0003-0653-4006 \\
Truman School of Government and Public Affairs and \\
Kinder Institute on Constitutional  Democracy, Smith College\\
and\\
Matthew Lavin ORCID:0000-0003-3867-9138\\
Data Analytics, Denison University\\
and\\
Richard Ball ORCID:0000-0003-4450-7957\\
Department of Economics, Haverford College\\
and\\
Sarah Supp ORCID:0000-0002-0072-029X\\
Data Analytics and Sustainability \& Environmental Studies, Denison University\\
and\\
Ben Gebre-Medhin ORCID:0000-0002-8140-0406\\
Department of Sociology and Anthropology, Mount Holyoke College\\
and\\
Laurie Tupper ORCID:0000-0002-4706-6205\\
Department of Mathematics and Statistics, Mount Holyoke College\\
and\\
Peter Kedron ORCID:0000-0002-1093-3416\\
Department of Geography, University of California, Santa Barbara\\
and\\
Nicholas J Horton ORCID:0000-0003-3332-4311\\
Department of Statistics, Amherst College\\
and\\
Joseph Holler ORCID:0000-0002-2381-2699\\
Department of Geography, Middlebury College
}
\maketitle
} \fi

\if0\anon
{
  \bigskip
  \bigskip
  \bigskip
  \begin{center}
    {\LARGE\bf Title}
\end{center}
  \medskip
} \fi

\bigskip
\begin{abstract}
Research communities are emphasizing reproducibility and replicability (R\&R) to improve the transparency and credibility of published research.
This transformational change to research practice requires a parallel transformation  in how students learn to conduct research in the first place.
Students need to start learning R\&R from the beginning of their undergraduate studies and continue to practice R\&R throughout their education. 
R\&R is fundamentally about maximizing transparency in planning and conducting research so that others can repeat the work and evaluate its claims. 
By integrating R\&R directly into foundational curriculum, rather than introducing it after teaching technical skills, students can develop replicable habits from the very beginning of their careers as knowledge producers. 


We identify four main principles of R\&R that can be integrated into existing curricula. 
First, R\&R is a habitual practice for conducting and engaging with research, not a property of research products to be checked at the eleventh hour. 
Second, the foundation of R\&R is intentional organization and clear communication of research materials. 
Third, R\&R begins with conceptualizing and articulating an analysis plan prior to analyzing data. 
Finally, R\&R is about critical reflection and evaluation of research by the researcher, peers, and broader community. 
We provide examples of how curricula can be transformed in a step-wise process that supports traditional learning goals with six classroom activities. 
Each of these activities introduces principles of R\&R at different stages of the research process as well as teaches students to be critical consumers of scientific research. 
\end{abstract}

\noindent%
{\it Keywords:} Reproducibility, Replicability, Pedagogy, Curriculum
\vfill

\newpage
\spacingset{1.8} % DON'T change the spacing!

\section{Introduction}

Principles and practices related to replicability and reproducibility (R\&R) should be integrated systematically and ubiquitously in quantitative methods training across disciplinary domains, beginning in introductory courses at the undergraduate level.
As standards for R\&R of quantitative research evolve and are strengthened, it is essential that new cohorts of scientists possess the knowledge and skills they will need to comply with new standards \citep{nasem2019}.
Even for students who do not go on to research careers, concepts and methods underlying R\&R reinforce principles of critical thinking and research integrity that are at the heart of college education and professional competencies.
Nevertheless, direct experience with the skills and practices that link R\&R to critical thinking are often absent from the classroom.
In this paper, we present four guiding principles and six classroom activities that introduce students to fundamental dimensions of R\&R.
These principles and activities demonstrate the feasibility of incorporating R\&R into the undergraduate curriculum and provide a resource for instructors. 

Although different disciplines have different interpretations of the terms reproducibility and replicability, we base ours on the interdisciplinary frameworks of the NASEM \citeyearpar{nasem2019} and Turing Way \citeyearpar{TheTuringWay}.
\emph{Reproducibility} is when a study can be independently repeated with the same data and procedures, producing the same results.
\emph{Replicability} is the ability of a study to be repeated with new data and the same procedures, producing similar results.
Replicability requires reproducible methods and also requires the research to be free of bias stemming from questionable research practices.
It follows that practices that support R\&R must ensure transparency in both the materials and methods of a study \textit{and} the original study purpose, design, and justifications for changes thereof.
The mechanics of R\&R are therefore well aligned with open science standards and practices \citep{nosek_promoting_2015}.

It follows that R\&R research practices must ensure transparency in both the materials and methods of a study and the original study purpose, design, and justifications for changes thereof.
However, higher education has left students largely unprepared to master reproducible research practices, or to navigate an expanding corpus of scientific literature with variable levels of transparency, and at times questionable research practices.
The scientific community has grappled with a replication crisis for several decades. Studies have been, and continue to be, published with minimal information about how results were obtained and conclusion were made. 
With numerous well-documented "failures to replicate" \citep{doyen2012behavioral, mobley2013survey, levelt2012flawed}, researchers have applied vigorous effort to establish and support better open science practices, many of which have focused on increased procedural transparency \citep{sandve2013ten, christensen2019transparent}. 
These innovations have yet to become a regular part of university education.

We need to transform higher education with clear integration of R\&R learning goals, starting from the beginning of undergraduate studies \citep{dogucu2024reproducibility}. 
In this scenario, students build a foundational understanding of how to be critical consumers of research and data, and are able to produce transparent and reproducible research.
Faculty should implement principles of open science when research skills are being learned across the curriculum.
Otherwise, students may develop habitual barriers to R\&R when they first learn research methods, working against any later introduction of open science practices and missing opportunities to engage more deeply with the research community. 
By integrating these procedures directly into the foundation of training researchers, we can instill norms of transparency into the research of the next generation of scholars. 
The \cite{NASEM:2018} report noted several components of data acumen that tie directly into R\&R initiatives.
Introducing these concepts at an early stage helps students foster good research habits that integrate principles of R\&R into the foundations of their research practices \citep{ostblom_opinionated_2022}, rather than later trying to change ingrained habits. 

Not only does teaching R\&R early and often integrate responsible research practices with learning research skills, it can also easily utilize best pedagogical practices such as active learning to reinforce course content and deepen student learning.
Preliminary evidence from the extensive and growing literature in open science suggests that if instructors create space to teach the motivational context and technical mechanics of R\&R, students benefit from gaining competency and self-efficacy to participate more fully in the academic community \citep{pownall_teaching_2023}.

The principles and activities highlighted in this paper apply to a wide variety of fields and methods, and at different stages of a research project, such as examining the work of others or formulating and implementing a research plan. 
We have designed activities that can be implemented for undergraduate students at any level, from first-year students learning about the principles of conducting research to upper-level students who are working on longer-term research projects. 

To address some potential concerns about the feasibility and desirability of incorporating R\&R in quantitative methods training, we use a question and answer format: 
\begin{itemize}
    \item Q: What if students lack the technical skills needed to learn R\&R? 
    \item A: It is becoming apparent that organized project management is a barrier to student success overall, which teaching replicability can help them overcome. Also, R\&R is about assessing the claims made by prior researchers. Learning to critically assess quantitative arguments can be taught as students develop technical proficiency. Moreover, developing this skill is a core component of developing quantitative literacy. The foundations of R\&R can be taught without prior knowledge of coding or requiring students to use particular technical tools.

    \item Q: Will replicability require too much time and effort to teach?
    \item A: It does take time, but practicing designing and executing clear and well-organized projects can be much more efficient than searching through permutations of analysis that did not start with a clear plan. In many cases, teaching replication can be integrated with activities that you are already doing in class, such that it doesn’t need to replace core content in the course.

    \item Q: Replicability is not consistent with my learning goals.
    \item A: Replicability concerns all research, regardless of discipline. Furthermore, replicability can enhance many existing learning goals by encouraging students to engage more deeply with any study or finding, critically consider project planning, thoroughly document analyses, and to clearly communicate the research process and resulting products.
    \item Q: Do my students really need to learn replicability?
    \item A:  Academic research is increasingly required to be replicable by publishers and funding agencies, and replicability is crucial in professional research environments. Furthermore, research has shown that training in replicability improves students’ ability to learn subject content from prior research, comprehension of their research projects, competency for independent and collaborative research, and confidence that they are able to participate in and contribute to the academic research community \citep{pownall_teaching_2023}.
\end{itemize}
This consideration of potential counter-arguments only strengthens the case that teaching principles and practices of R\&R in quantitative methods courses should be the norm.
We propose four principles for R\&R that should guide teaching and learning. 
First, reproducibility should be viewed as an ongoing practice that informs how research is conducted and shared rather than as a static property of completed projects.  
Second, the foundation of transparency and R\&R is intentional organization of research materials, both while work on a project is underway and when the documentation for a completed study is assembled.
Third, committing to pre-analysis constraints, such as pre-registration, enhances research integrity by reducing biases and questionable research practices.
Fourth, R\&R is about critical reflection and evaluation of research by the researcher, peers, and broader community.

Building on these principles, we introduce a series of activities designed to help students internalize the practical and ethical dimensions of R\&R.
These activities include reviewing the reproducibility of current research, creating reproducible projects and documentation, sharing research projects with peers, drafting analysis plans, and conducting peer evaluations to reinforce reproducibility as a collective responsibility.
They are designed for use in diverse disciplines and at all levels of the curriculum, from introductory courses to research seminars. 

By embedding reproducibility into the curriculum, we aim to foster a culture of transparency and accountability that extends beyond individual projects.
Additionally, for courses not focused on training technical research skills, integrating R\&R principles can make students better knowledge consumers.
This paper contributes to the broader conversation on research integrity by offering educators practical, adaptable strategies to cultivate a ``reproducibility first" mindset among students.
Our approach highlights the interplay between technical execution and ethical responsibility, ensuring that students not only learn how to make their research reproducible, but also appreciate its broader significance within the scientific community. 


\section{Guiding Principles} \label{sec-principles}
We present four cross-cutting principles to educators teaching replicability. The following principles can be used when forming learning goals and the pedagogical design of courses. 
These four principles cross-cut each of our learning activities and inform our learning goals and our approach to pedagogical design.
Our principles are as follows:

\begin{enumerate}
    \item Reproducibility is a \textbf{practice}, not a property.
    \item Intentional \textbf{organization} of research materials is the foundation of transparency and reproducibility.
    \item ``Tie yourself to the mast'': Reproducible research establishes pre-analysis \textbf{planning} constraints.
    \item R\&R makes a commitment to improving knowledge through critical \textbf{reflection}.
\end{enumerate}

\noindent The example activities in Section \ref{sec-activities} aim to synthesize these principles, with each activity having its own unique emphasis; some focus more on building technical skills and others more on higher-level ethical considerations.
Each principle is explained in more detail below. 

\subsection{Reproducibility is a \textbf{practice}, not a property}

Replicability is an ongoing practice, not just a property of published research or a requirement to fulfill at the eleventh hour of submission.
Replicability should not be understood only as a property of published research projects: it should drastically change practices in the way students learn to conduct research and engage with the research community. 
Replicability is often put into practice and trivialized as task list of reproducibility requirements for publishing and as property of published research: are the data available and does the code run and produce the results?
We argue that replicability is better understood and taught as a set of practices for producing and engaging with research.
The practice of replicable research should begin from the moment that we form a research question.
Replicability enacted in this way adds value to our process of learning and creating new knowledge. Instead of asking if a published work checks the right boxes, it is meaningful to ask how we might best learn from prior studies to solve problems in new contexts.
In evaluating research by others, we should ask whether the authors have made it possible (and ideally easy) to reproduce, validate, learn from, and extend their work.
The activity described in section \ref{sec-recognize} provides an opportunity for students to evaluate research by others.
In our own research, we should in turn ask how we can best proceed so that others can easily reproduce, replicate, validate, learn from, and extend our research.
This reorientation of replicability from a property to a practice implies that students shift from a paradigm of post-research evaluation to learning how to participate in the scientific community.


\subsection{ Intentional \textbf{organization} of research materials is the foundation of transparency and reproducibility} \label{proj-org}

Transparency is essential to the critical assessment of researcher claims through the reproduction and replication of research, and intentional organization is the key to transparency. 
This is also true for students who complete a class project to demonstrate their knowledge or apply their research skills.
Project organization is built around establishing and articulating a well-defined structure of work. 
In quantitative and computational studies organization often centers on the development of a directory structure–a well-defined hierarchy of folders and subfolders–that stores the data, scripts, and other files used during a study. 
Researchers should build the directory structure they will use for a project before they start working with their data, and then populate the folders with the various files–data, scripts, output, and supporting documentation–they use throughout their work on the project. 
In classes, we can require students to do the same, while teaching them the reasons why careful organization is critical to strong argumentation and trust in evidence. 

Most obviously, building a directory structure at the beginning of a project helps researchers “keep track of their stuff”. 
Each time a new file is obtained or constructed, they know where to store it, and when they need to locate an existing file, they know where to find it. 
Instructors should not take it for granted that students will find it natural to organize their work in this way. Anecdotal accounts from faculty suggest that instructors must introduce these concepts to students, who are often accustomed to applications and search engines that no longer require an understanding of directory structure (Chin 2021). 
Without guidance, students may not pay attention to where files they download or create are stored, and they may have difficulty finding files when they need them. 

Intentional project organization is critical to reproducibility in ways that go beyond simply keeping track of your stuff. 
When you have completed a project, the directory structure you built at the outset, containing the data, scripts, and other files you assembled throughout the course of your work, can serve as reproduction documentation. 
Using this documentation, any interested third party should be able to reproduce all the computations you executed and all the results you reported simply by running the scripts. 
Such “automated reproducibility” is possible, however, only if the documentation you have assembled is organized in a way that satisfies certain principles. These principles are summarized in the \emph{Reproducibility Trifecta} \citep{ball_2023}:

\begin{itemize}
    \item \textit{Define a directory structure: }All the documentation for a project should be stored in a well-defined hierarchy of folders and sub-folders.
    \item \textit{Make explicit choices about the working directory: }For every script that is written, a deliberate decision must be made about what folder should be designated as the working directory when the script is executed, and that decision must be communicated in the documentation (for instance, in a {\tt Readme} file, or in comments in scripts).
    \item \textit{Use relative file paths: }In all scripts, whenever it is necessary to identify a certain folder (e.g., one containing a file that needs to be opened, or one in which a newly created file should be saved), the location of the folder should be specified using a file path that begins in the current working directory (not in the root directory of any particular computer).
\end{itemize}

These three elements of the Reproducibility Trifecta are interrelated and need to be coordinated.
When you write relative file paths in a script, you must know which folder will be designated as the working directory, because that is where the paths should start.
You need a well-defined directory structure because you must specify a valid path through the hierarchy of folders and sub-folders containing your documentation.
Organizing a project in a way that achieves this harmonization of directory structure, working directory, and relative file paths is an essential foundation of reproducibility.

The activity described in section \ref{sec-organize} introduces students to the principles of project organization discussed above.
It consists of an exercise in which students do some simple steps of data wrangling and then construct basic numerical and graphical summaries.
The focus of the exercise is on the practices that ensure reproducibility---assembling comprehensive documentation that satisfies the principles of the Reproducibility Trifecta. 

Not only do students need to know how to organize their own data, they need to know how to organize it in a way that allows for the data and code to be easily shared with others.
This includes standardizing file naming conventions, using relative file paths, and documenting any software used to clean data or run an analysis.
Clear organization and detailed documentation reduce barriers to collaboration and replication.
By building a workflow that can easily be shared or uploaded, students learn a critical part of constructing a reproducible project.

The activity described in section \ref{sec-rotation} shows students the importance of understanding how one's project is organized.
By having students share their data and code with one another, they must build the project in a way that makes it reproducible.
File paths must be relative, proper software must be loaded, and students must know where their data is in order to share it.
Additionally, by receiving a project from another student, they learn the importance of contextualizing the analysis through comments describing decisions.

Both of these activities take a hands-on approach to make students consider the structure of their workflow and organization of their data.
Developing a technical understanding of file paths, workflows, and reproducible coding practices lays a foundation for students to engage in reproducible research.

\subsection{``Tie yourself to the mast'': Reproducible research establishes pre-analysis \textbf{planning} constraints} \label{sec-mast}

``Tie yourself to the mast": Reproducible research establishes and holds to pre-analysis constraints, so build a ``mast" to guide your research and tie yourself to it.
Research is designed and implemented in a complex social world full of subjective researcher decisions with a range of possible outcomes.
Where the first principle focuses on the technical knowledge needed to engage in reproducible research, our second principle emphasizes engaging in ethical practices to make research more replicable.
More specifically, we see a need to teach students the importance of how to ask research questions and plan analyses before the results are known.
Clearly defining the questions you are asking before starting the analysis can help mitigate the influence of uncertainty and subjectivity in one's work, and especially the temptation to work backwards from a surprising result as if it were always present in a researcher's hypothesis.
Developing research questions and designing research methods based on literature review---before conducting any analyses---is consistent with principles of transparency and accountability in the research process.
Analysis plans should address major sources of bias, which can be specific to the discipline and type of research \cite{nosek_preregistration_2018}.
Generic and discipline-specific analysis plan templates prompt researchers to think about and declare their research design decisions in advance, thus reducing bias and potentially improving the replicability of empirical research. 

We engage with this principle in several ways.
The activity described in section \ref{sec-preregistration} asks students to pre-register their project before analyzing models.
Whether they conduct the analysis or not, this activity emphasizes the importance of developing clear, testable hypotheses with appropriate research designs before knowing the results.
Pre-registration helps students internalize the value of committing to a research plan, reducing the temptation to alter hypotheses post hoc.
This activity gives students first-hand experience developing a research plan, and then sticking to it regardless of whether the results are statistically significant.
By making a plan, and sticking to it, students can understand the importance of developing a transparent research design.

With respect to p-hacking, the activity described in section \ref{sec-dubious} emphasizes how to recognize dubious practices that include overfitting models, making errors in extrapolation, and ignoring or misinterpreting null results.
A point that we want to reinforce with students is not that there is a single way to conduct an analysis, but that researchers must be transparent in their decision-making process and acknowledge the limitations behind their research.
Cultivating this transparency equips students to engage in ethical, replicable research.
These activities provide a microcosm of the research process for students and provide practical experience in modeling and motivating ideal research processes and communities in the classroom. 

\subsection{R\&R makes a commitment to improving knowledge through critical \textbf{reflection}}

In an open science framework, the fundamental purpose of R\&R is to facilitate self-correction in science \citep{nosek_promoting_2015} with the overarching goal of improving knowledge. 
As members of a research \textit{community}, scientists are obliged to verify and correct errors in the published literature through review and replication. 
In turn, as members of a classes or cohorts in higher education, students can be trained and encouraged to constructively review each others' work. 
As individuals, students can improve their own work and learning through self-reflection. 
Critical reflection is a cornerstone of R\&R, but committing to R\&R also entails making a commitment to the importance of critical reflection on knowledge improvement. Self- and peer-reflection happens at multiple levels (self, class/cohort, and community).
What makes reflection critical is not negativity or criticism for its own sake. 
Effective reflection is not merely cynicism. 
In contrast, the central aim is to improve research quality overall \citep{janz_replicate_2021}, which involves identifying flawed approaches (so that they might be improved upon) and identifying best practices (so that they might be extended and duplicated).

Lastly, critical reflection finds a natural ally in active learning pedagogy. 
Reflective assessments (self-assessment, peer-assessment, etc.) can be viewed as modalities within the broader category of active learning, in that they move away from education as teacher-to-student knowledge transfer by asking students to ``think, reflect and be curious'' \citep{tuta:2022}. 
Active learning strategies are generally thought to increase students' intrinsic motivation and to develop their capacity for meta-cognition.
Meta-cognition, or thinking about one's thinking, is a specific type of critical reflective that has been linked to numerous positive education outcomes \citep{dign:2008,fleu:2021}.



\section{Sample Classroom Activities} \label{sec-activities}

We now propose a series of classroom activities designed to bring the principles outlined above into practice.
These activities serve as a starting point for instructors and are flexible enough to be applied at various levels across a wide range of disciplines and methodological approaches.
While most activities involve the use of a scripting or programming language, many can be easily adapted to different classroom environments and skill sets.
Detailed instructions for all activities are available on a publicly accessible repository, which can be found \href{https://github.com/lacombe129/R_R_vignettes}{here}. The activities range in scope from ones that may be implemented in a single class to semester-long, scaffolded projects. We see opportunities to extend on each of these activities for more advanced students or research-intensive courses, and provide potential ideas for activities to further implement the principles we highlighted above. 

The activities are structured to build toward a deeper understanding of the four core principles. 
They begin with an emphasis on foundational concepts, followed by the development of technical skills necessary to ensure reproducible data cleaning and analysis, and culminate in fostering critical thinking skills to evaluate both others' work and students' own research.
The first three activities primarily focus on Principle 1.
Activity 1 provides guidance on how students can identify and assess the implementation of reproducible and replicable research practices in published work.
Activities 2 and 3 focus on embedding reproducibility at the outset of research, first by organizing research materials in a reproducible manner, and then by reinforcing reproducibility through collaboration with other groups in a rotation, ensuring that each group builds upon the work of others.
Activity 4 addresses Principle 2, encouraging students to plan their analysis before beginning data work—a shift that integrates replicability principles from the start of the research process.
Finally, Activities 5 and 6 emphasize the importance of critical thinking skills highlighted in Principle 3, helping students assess reproducibility as an ongoing practice to be woven throughout the research process, both in their own work and when evaluating others.

\subsection{Recognizing Reproducible and Replicable Research} \label{sec-recognize}

\textbf{Learning Goals and Outcomes:} 
In order to achieve the aspirations of reproducibility to enhance transparency and improve scientific communication, researchers and students alike need to learn how to read ``reproducible'' research.
For the purposes of reproducibility, ``reading'' must also include understanding the supplementary data and procedures associated with research and ability to execute any code or computational models \citep{nasem2019}.
Acknowledging that the reproducibility of any research project cannot be assumed, students will also need to learn how to identify all of the components required for reproducibility and to assess and test their completeness and functionality.

Through this assignment, students will produce a reproducibility review of a prior study analogous to published literature reviews (e.g., \cite{Hardwicke2020,Collins2022}) assessing the reproducibility of particular domains of research.
These efforts are not trivial and need to be learned \citep{vilh:2022}.
The learning goal of this activity emphasizes the external contextual motivation for reproduciblity by guiding students through an evaluation of prior published research.
Students will gain awareness of the required components of a reproducible study and appreciation for the difficulties and importance of improving reproducibility in their discipline.

\textbf{Scaffolding/skills:} This assignment should follow a conceptual introduction to scientific reproducibility and shall build on students' skills in reading scientific research papers, searching for and evaluating data sources, reading and implementing research procedures in a particular domain of science, and scientific communication.
The assignment scaffolds on the concepts of reproducibility with the active learning exercise of auditing published research.
Faculty may reinforce concepts by directing students to evaluate studies in a field before and after reproducibility standards were put into place, or to evaluate an original version of a study published without reproducibility standards and a later reproduction of that study published with those standards. 

Students should start the assignment having familiarity with the structure of empirical research papers and experience reading and understanding their contents.
Students will likely need to read and re-read papers more closely as they scrutinize the narrative for procedural details.
Furthermore, they will search for and scrutinize supplementary materials and references to data and procedures to gain a more holistic understanding of the study. 
Students should have some background in research design and procedures in their field.
Ideally, they will assess prior studies which reinforce knowledge of methods in their field, and gain experience reading about empirical materials and methods in their field.
Finally, students will build upon their research communication skills by following a protocol to audit reproducibility and reporting their results according to the protocol.

\textbf{Implementation:} We recommended selecting the study or studies for review to maximize learning outcomes in the context of the curriculum.
Studies could be selected based on compatibility with the method(s) and topic(s) in the curriculum, and as candidates for future reproduction and replication research.
If students are to select studies on their own, the literature search parameters should be constrained to support learning outcomes.
The search can be constrained to empirical studies that are also published in selected journals and use particular software environments (e.g. R or Python), methods, (e.g. linear models), or data sources (e.g. the Census). 

Once studies are selected, the review activity can be implemented by following a relevant reproducibility assessment protocol or customizing a protocol for their needs.
For example, the Social Science Reproduction Platform covers reproducibility reviews in chapters 2 and 3 \citep{BITSS2020} and the Association for Geographic Information Laboratories in Europe (AGILE) conferences have a two-page guide for peer reviewing the reproducibility of conference papers \citep{nustAGILEguidelines}.

Effective protocols should describe the discipline-specific standards for a reproducible research paper and guide students through reading the research materials and reporting findings.
Reading guidance should focus on how much effort students should reasonably expend.
For example, AGILE reviews \citep{nustAGILEguidelines} list the ``do's and dont's'' of reviews, including ``Do look for a README file with step-by-step instructions, and do not dig through poorly documented files to figure out which files produce outcomes in the paper."
Reporting guidance is essentially a rubric for describing the reproducibility of a prior study in the form of checklist items, ratings, and/or a narrative report. 
Examples of such outputs include the supplemental materials of \citet{vilh:2022}, the AGILE reproducibility reports, and the AGILE reproducibility reviews \citep{AGILE2024}.

This activity may be abridged or expanded to meet constraints and learning goals.
It may be simplified by providing a specific prior research study and step-by-step activities and questions to emphasize particular reproducibility concepts.
For example, instructors could emphasize the importance of intellectual property licenses by asking students to look for license or metadata files describing the legality to copy each dataset used in a study. 
An abridged version could be framed as a scavenger hunt for research artifacts in a manuscript and its supplemental resources.

Example scavenger hunt questions:
\begin{itemize}
\item Which version of Python was used for the study?
\item Which packages are used for the study?
\item Which code files can be used to run the study, and in what order?
\item Describe the data sources used for the study and find a link to download them.
\end{itemize}

The activity may be expanded by going beyond the question of whether particular reproducibility components exist into a more thorough investigation of how complete the components are and how closely the methods and results as described in the research paper. 
For example, if a study uses Python code to compute and visualize a local Moran's I test for spatial autocorrelation, then students should draw out graphic workflows of the methods as described in the paper and compare line by line with a close reading of the code to determine if the code accurately and completely matches the manuscript. 
If a study states that it is reproducible because it uses publicly available data sources, then students could attempt to verify the availability of those data sources and that they contain the same variables and number of records as described in the manuscript.
The most thorough audits will also include checking completeness and accuracy of data and procedures \textit{vis a vis} the study as written, and review of the study design's conceptual and internal validity.

\textbf{Assessment:} The learning goals can be assessed through evaluation of the completeness, accuracy, and clarity of the students' reproducibility assessment. 
In the abridged version of the activity, the instructor will already know the true status of the prior study and can simply evaluate the students for completeness and accuracy.
If a sufficient number of individuals or groups assess the same study, their findings can be compared to one another.
If the class continues to investigate the study, results can be compared against further discoveries through reproducing and replicating. 

In the expanded version of the activity, instructors can evaluate student's articulation and application of concepts related to research practices and validity. 
Finally, to reinforce learning objectives, students should be prompted to write reflections or hold in-class debates on their findings in the context of reproducible science. 
What are the implications of their findings for both the study at hand, and for knowledge production in their field? 

\textbf{Possible Extensions:}
This exercise can naturally be extended with the activities to come.
Building on section 3.2, students may reassemble the research project components according to specified organizational structure.
Building on section 3.4, they may be tasked to write a pre-analysis plan for reproducing the study or replicating the study in a new context. 
Building on section 3.5, they may research the assumptions and best practices for the research design used, and use the knowledge gained from this exercise to further assess not only the study's reproducibility, but also its validity. 
Integrating section 3.6, students may synthesize each others' findings to form the basis for a broader reflective discussion about reproducibility in their field.
With all of these foundational principles and practices in place, students could even take on the task of attempting to reproduce the paper's results on their own, but we caution that this is a challenging task \citep{vilh:2022} that will require multiple weeks of effort.

\subsection{Project Organization} \label{sec-organize}

\textbf{Learning Goals and Outcomes:} In this activity, students gain hands-on experience with research methods that embody the concepts of project organization discussed in section \ref{proj-org}.  

The activity consists of an exercise in which students use data from a cross-section of countries to investigate whether there is any association between income per capita and the average level of well-being (as reported by respondents to a socio-economic survey).  
Students combine and clean data from two sources, construct tabular and visual summaries of the processed data, and write a brief report in which they answer several questions about their findings. 
In addition to this report, students submit documentation that can be used to reproduce all the computations they carried out for the exercise.  
The instructions given to students, as well as an example of the report and reproduction documentation a student would submit, are available on the GitHub repository.

The data processing and analysis involved in the exercise are intentionally very simple; the focus is on methods of conducting and documenting work with quantitative data in ways that ensure a high level of transparency and reproducibility.
The conventions of project organization adopted in this exercise closely follow the Teaching Integrity in Empirical Research (TIER) Protocol \citep{tier}, but instructors should feel free to modify the organization of the exercise in any ways that suit their purposes.
The ultimate purpose is to equip students with skills in reproducible methods that they can apply in future research projects, including papers for other classes, theses, or dissertations.


\textbf{Scaffolding/skills:} Before beginning this exercise, students should have at least a basic understanding of how to use scripts or batch files to execute computations, as opposed to using point-and-click menus or interactively typing one command at a time.
It is also helpful if students are familiar with the concepts of the working directory and relative directory paths, but if students lack that background a simple lesson or demo can be sufficient to bring them up to speed before embarking on the exercise.  

The exercise is highly scaffolded, in the sense that the instructions explain the tasks students must carry out in a sequence of discrete steps.  
Moreover, the instructions are written in simple, explicit language that should not be difficult for students to understand, even if they have little previous computational experience.

The key principles of project organization, in particular the Reproducibility Trifecta, are introduced incrementally as students progress through the exercise.  
The first task is to establish a directory structure where the files for the exercise will be stored, with the instructions specifying what that structure should be. 
The instructions also prescribe a particular convention for managing the working directory, namely that whenever a script is being executed, the top-level folder in the directory structure should be designated as the working directory.  
The instructions for writing the scripts that process the data and generate the results remind students to use relative directory paths.

Although the key elements of reproducibility are introduced in small steps, they generate a large payoff:  students learn to create documentation for a project that would enable a third party to independently reproduce their work. 
When they have completed the exercise, students submit not only the written report, but also the reproduction documentation they have assembled.
If the project is completed successfully, the instructor should be able to generate the results presented in the report simply by executing the scripts.

\textbf{Implementation:} If possible, we recommend distributing the instructions and giving students a brief overview of what the exercise entails at the beginning of a class or lab session, and then letting students spend the rest of the period beginning their work on it.
If students make a good start during that initial class or lab, they should be able to complete the exercise outside of class time.

The first step in the exercise is constructing the directory structure in which students will store their work for the exercise.  
This task is straightforward, and explained in detail in the instructions, but it is still important to stress to students that this is essential and they must actually do it before proceeding.  
Students are often (mistakenly) view the task of building a hierarchy of empty folders and sub-folders as trivial or vacuous, and are therefore tempted to skip it.  
In addition, when students do create all the necessary folders, it is important to be sure they really know where they are.  
In particular, they may not understand whether they are working locally on their personal computers or on a cloud platform like Microsoft OneDrive or Apple iCloud, and that can lead to confusion.

After building the directory structure, students write two scripts: one that processes the data from the two input data files, and one that uses the processed data to generate the results that students present in their reports.  
Automated reproducibility depends on the harmonization of the relative file paths used in those scripts with the directory structure and the convention chosen for managing the working directory.
Again, the instructions give detailed guidance about how to achieve this harmonization, but in light of the crucial role this plays in achieving reproducibility, some reinforcement from the instructor can be valuable.

When the exercise is complete, instructors may choose whether to ask students to turn in their reports in printed or electronic format.  
To turn in their documentation, students zip the top-level folder of their directory structure, and submit it electronically.

\textbf{Assessment:} The fact that students submit reproduction documentation for this exercise has little if any effect on how the report is assessed.
The instructor may take whatever approach to reading and commenting on student work they are accustomed to.

In assessing the documentation, the instructor first considers whether the project is organized as specified in the instructions, and whether the steps required to reproduce the student's work--in particular, the order in which the scripts should be run--would be clear to a third party. 
The main task is then to try executing the scripts, to see whether they run without any errors and produce results that match those presented in the report.  
When it is successful, performing these checks should be quick and easy.  
Finding out what the problem is when a script does not run is an additional step, but because the instructor has access to the scripts, it is usually not difficult to pinpoint the issue.

Exercises like this one, in which students construct reproduction documentation for a project, also lend themselves well to peer assessment.  
Working in pairs, students can try executing each other’s scripts and seeing if they get the intended results; if they don’t they can work together to find the problem and fix it.  
This cooperative exercise can enhance student learning, and also reduce the amount of work the instructor must do for the assessment.

\textbf{Possible Extensions:}
The exercise presented in this section provides a model of how an instructor can transform any existing homework problem or lab exercise into a reproducible project by providing guidance on the three elements of the reproducibility trifecta---defining and building a directory structure, being attentive to the working directory, and using relative file paths.  
In addition, advisors of theses, dissertations, and other independent research can provide guidance to their students on how following these principles can ensure that their work is transparent and reproducible.

Incorporating reproducibility extensively in quantitative methods instruction is not only possible, but should be standard practice.  
Students should begin learning principles of project organization and documentation when they first start working with statistical data, and continue learning and practicing reproducible research methods through every phase of their training.

Our vision is that once reproducibility is ubiquitously infused throughout the curriculum, using reproducible methods will become the norm among practicing researchers.


\subsection{Rotation-Based Classroom Project} \label{sec-rotation}

\textbf{Learning Goals:} Once students understand how data and files are organized on their computers, the next challenge is learning how to structure a project to support collaboration and reproducibility.
This includes not only organizing files and using relative file paths, but also managing software dependencies—such as specifying package requirements, identifying the version of the software being used, and indicating which data files must be loaded.
By sharing a data project with a collaborator, students gain valuable insight into the full set of dependencies and contextual information (e.g., comments and documentation) necessary for another researcher to successfully run their code.
When they send the project to another student, they must consider what is needed for a project to run.
Likewise, when receiving a project from a peer, students have the opportunity to critically assess what elements support reproducibility and what improvements could be made to enhance clarity and functionality.


\textbf{Scaffolding/skills:} To address these issues, we built a classroom project as part of a one-credit hour pass/fail lab introducing the fundamentals of R using Tidyverse.
This course first introduced concepts such as how to load data (both on local directories and remotely), produce polished figures and plots, and how to clean data. 
After learning the basics of R and Quarto, students were tasked with identifying a potential dataset to use for future iterations of the course and writing a report that describes variables in the dataset, producing visualizations and summary statistics of key variables, and discussing the ethics of how the data were produced.
The final report is roughly two to three pages and produced in Quarto.
Rather than making the final project a separate assignment, it was integrated into weekly labs, so students iteratively completed each component of the report over the second half of the semester. 
This helped reinforce weekly course material and also helped students understand that conducting research is a step-by-step process, with space for reflection and review.

\textbf{Implementation:} Students form groups of three to four, each with their own unique proposed dataset.
After the groups are established and the data is briefly introduced, students then `rotate' the project by sending both the data and the .qmd (Quarto) file to the next student in line to complete the next task for the report (see Table \ref{table:rotations} for an example of the rotations).
By rotating the project at each stage, students know that they need to write code in a way that will run on another computer.
File paths must be correctly identified, and all required packages must be included in the code.
Over the course of five weeks, the students continue a weekly rotation, moving from introducing the data, to selecting key variables, summarizing these variables, and visualizing them. The final rotation is focused on data ethics. 
The final rotation is focused on data ethics.

In addition to this format being helpful in introducing principles of reproducibility to students, it also was designed so that every student had to complete every part of the project.
Students could not divide the work to play to their strengths.
The students with a stronger quantitative background still needed to write, and those with stronger writing skills still needed to complete their assigned tasks in Quarto individually.
The rotations facilitated each student's need to demonstrate competency for each part of the project.
At the end of the project, students complete a self-reflection to evaluate both the process of collaborating on a project with others, as well as some of the ethical considerations around data provenance and how the chosen datasets can be used. 



\begin{table}[ht]
\resizebox{\textwidth}{!}{%
\begin{tabular}{lccc}
\multicolumn{1}{c}{}                                  & Data   Option 1                  & Data   Option 2                  & \multicolumn{1}{l}{Data   Option 3} \\ \hline
\multicolumn{1}{l|}{Rotation   1: Data Proposal}      & \multicolumn{1}{c|}{Student   1} & \multicolumn{1}{c|}{Student   2} & Student   3                         \\
\multicolumn{1}{l|}{Rotation   2: Variable Selection} & \multicolumn{1}{c|}{Student   2} & \multicolumn{1}{c|}{Student   3} & Student   1                         \\
\multicolumn{1}{l|}{Rotation   3: Summary Stats}      & \multicolumn{1}{c|}{Student   3} & \multicolumn{1}{c|}{Student   1} & Student   2                         \\
\multicolumn{1}{l|}{Rotation   4: Data Viz}           & \multicolumn{1}{c|}{Student   1} & \multicolumn{1}{c|}{Student   2} & Student   3                         \\
\multicolumn{1}{l|}{Rotation   5: Data Ethics}        & \multicolumn{1}{c|}{Student   2} & \multicolumn{1}{c|}{Student   3} & Student   1                        
\end{tabular}%
}
\caption{Sample Project Rotations for a Group of Three Students}
\label{table:rotations}
\end{table}

\textbf{Assessment:} After all rotations are completed, students choose their preferred report to submit and work together to finalize formatting and include all necessary components. The instructor then assesses these projects on a pass/fail basis. 
Projects are assessed to ensure they have all necessary components, that the code can run and the report can be rendered by the instructor, and that the report is formatted according to project guidelines (captions on all images, all code hidden, etc.). 
Students also individually complete a self-assessment where they reflect on the process of sharing their code and data, including the challenges their groups faced. 

Assessment for this project was based on two components: submission of the group’s chosen report, and peer assessment. The peer assessment component of the final project is an opportunity for students to discuss how the rotations went. Rotations themselves are not assessed by faculty week to week due to the number of students in the introductory courses (30 per section, with faculty teaching on average 3 sections), so the peer assessment allows students to report any concerns or add context to the group dynamics. If group members felt that every team member completed the rotations, then all members earned a token for submission. 


\textbf{Possible Extensions:}
This project was designed to introduce students to the fundamentals of using Quarto, but the general structure of a rotation-based project could be extended to almost any course, either quantitative or qualitative. 
For courses using data analysis, this project could work for most software, including Stata, Python, etc. 
This activity can also work in a qualitative methods course.
For example, students in four-step rotation could conduct and audio-record interviews, transcribe them, develop a qualitative coding scheme, and analyze results.

\subsection{Pre-Registration/Plan of Analysis} \label{sec-preregistration}

\textbf{Learning Goals and Outcomes:} Students who are working on empirical research projects often express a strong desire to find significant results that support their hypotheses.
This pressure stems from a worry that null findings may be uninteresting or unimportant, or that findings counter to their expectations indicate some inadequacy in the report or methodology, and pre-registration addresses two primary problems that might emerge in the face of this pressure.
First, students may look to change the model specification or sample, modifying the analysis until significant results are found.
Second, if the findings do not comport with the original hypotheses, students may change their theoretical expectations to comport with their findings. 
However, one learning outcome we want to emphasize is that knowledge generation is difficult, and if we already knew the answers to our research questions, we would not be asking them. 
If our goal is to build knowledge, that means we need to accept that our theories may be wrong, and our findings may not be statistically significant, or lead us to unexpected conclusions. 
More specifically, we want to avoid hypothesizing after the results are known (HARKing) \citep{kerr1998harking}, or changing our expectations or analysis simply because the findings fail to support our hypotheses.

One approach to emphasize principles of reproducibility and replication is to incorporate pre-registration into the scaffolding for a project \citep{king2016registered}. 
Pre-registration has emerged as one possible solution to minimize the issues such as p-hacking, cherry picking data, or post-hoc hypothesizing. 
Before students have analyzed or collected data, they must formalize their expectations and outline their research design for instructors to review. 
This framework can be easily integrated into existing projects in courses using mixed methods or qualitative methods. 
Alternatively making a pre-registered plan of analysis can also be a project in itself. 
This approach can be particularly useful when time constraints make including analysis difficult for a semester-long course, or students do not yet have the required methods training to collect, clean, and analyze findings. 
This approach also centers the importance of grounding research in existing literature and developing a robust theoretical framework before analyzing or even exploring data. 


\textbf{Scaffolding/skills:} Pre-registration can be incorporated into any assignment or project that requires students to conduct data analysis. 
For a lab assignment, students can be assigned pre-lab work where they state their hypotheses or outline how they will be conducting analysis (either for qualitative or quantitative work) before they are given access to the data. 
This could be incorporated into a single class session where faculty give students material to develop their own expectations/hypotheses. 
For larger projects, pre-registration can be integrated into project deadlines scaffolded throughout a course. 
An advantage of this approach is that each deadline/check-in can be paired with small lectures or readings highlighting the purpose of each piece of a research paper (literature reviews, hypotheses, research design, etc) as well as how students should be thinking about incorporating it. 
Students can first develop their theoretical expectations and submit pre-registered hypotheses with their literature review or theory sections. 
After expectations are set students can then think critically about how best to operationalize and evaluate their hypotheses. 
One of the most appealing elements of this activity is that it is highly flexible in scope and technical training for students. 
For introductory classes, this assignment can be paired with discussions of different types of modeling and the choices researchers must make when deciding and justifying decisions to operationalize key concepts. 
For other courses where students may not be estimating models or analyzing quantitative data, students can use this process to highlight what empirical support for their hypotheses might look like. 

\textbf{Assessment:}
In addition to reinforcing principles of R\&R, this approach also provides an opportunity for faculty to give feedback on students' general plan of inquiry before analysis begins. 
Faculty can assess the plans for feasibility and appropriateness, and flag any areas of concern for students before they begin analyzing their research question. 
Plans can be evaluated for the strength of their theoretical development, appropriateness of model specification, and feasibility of the projects. 
This provides an opportunity to course-correct projects where needed before data collection has begun, and permits students to follow through with analyses even when results are not significant.

\textbf{Possible Extensions:}
The scope of the plan of analysis can grow or shrink to almost any size depending on the nature of the project. 
For students with more advanced technical skills, they could be assigned already published pre-registered studies or plans of analysis, and then implement the plan themselves using existing data or replication materials. 
This provides students an opportunity to use their technical skills using real, published work, while also allowing them to evaluate how well the plans articulate their plans of analyses, and identify ways it can be strengthened. 
For less advanced students, groups could be assigned published, pre-registered plans of analysis, and then evaluate the final published paper to see the extent to which the plan of analysis was followed, and whether deviations were explained.  

\subsection{Appraising and Assessing Dubious Research Practices} \label{sec-dubious}




Table \ref{table:transparency} describes a way to separate out issues of transparency/reproducibility and rigor.

\textbf{Learning Goals and Outcomes:}
It's important that students can identify, compare, and contrast dubious practices that may be hidden in a researcher’s methods.
This includes being able to make sense of a world of nulls and p-values, the impact of large samples, and the effect of multiple tests and/or high dimensionality.
To help students be able to assess the potential impacts of dubious practices, they need practice in settings where dubious practices are introduced.
Dubious practices in statistical modeling include Type I and Type II errors, conflating different measures of performance (significance, predictive performance, effect size/coefficient), over-fitting of models, and extrapolation error.
Students can deepen their understanding of reproducibility practices by generating spurious conclusions from data in settings where the truth is known, then scrutinize and reflect upon these approaches.
In geographic modeling, they could include modifying the geographic extent, scale, data classifications, and variable normalization to achieve the expected result or obscure undesirable patterns.

\begin{table}[!ht]
    \centering
    \begin{tabular}{llll}
    \hline
        Type & Transparent & Rigorous & Description \\  \hline
        1 & Yes & Yes & Gold standard of scholarship \\ \hline
        2 & No & Yes & Must take scholar’s word for it \\ \hline
        3 & Yes & No & Can be critiqued and improved upon \\ \hline
        4 & No & No & Hiding among type 2; unknown quantity \\ \hline
    \end{tabular}
    \caption{Cross-classification of transparency and rigor.}
    \label{table:transparency}
\end{table}

\textbf{Implementation:}
Much of the emphasis of the previous activities have been to promote principled practice to foster reproducibility.
However, it is also important to have students engage with practices, methods, and publications that are less principled, as a way for them to start the gauge the potential impact of non-reproducible and non-replicable approaches to scientific research.

Students need to be warned about the perils of HARKing, the file drawer problem, P-hacking, selection bias, and so on.
This activity is intended to help them see the impact of such practices in a specific example.
The class activity would allow students to explore the implications where a statistical procedure is used that doesn't maintain appropriate Family Wise Error Rate.
(Such an activity would help reinforce the importance of pre-registration of a protocol, see sections \ref{sec-mast} and \ref{sec-preregistration}.)

Teaching students how to ``creatively game" the prior research literature has a potential downside: students may become distrustful of all claims in statistical research. 
Our goal for this exploration is to balance preparing students to deal with a world where reproducibility guidelines are not uniformly practiced without leading them to conclude that all research is fatally flawed \citep{edmans}.

The perils of multiple comparisons when analyses have not been prespecified have been highlighted previously.
\cite{austin:2006} reported the potential issues when testing pairwise differences in astrological (Zodiac) signs without accounting for multiple comparisons when assessing health differences.
Students an implement this activity using general purpose statistical software and datasets with date of birth and multiple health outcomes.
Here there are 66 possible pairwise comparisons that can be made between all of the 12 Zodiac signs (e.g., Capricorn vs.\ Leo, Aries vs.\ Taurus).
If all of the outcomes were the same for each of the Zodiac signs in the population, we would expect 5\% of the comparisons to yield a p-value $< 0.05$.
However, the probability of one or more of these pairwise tests yielding a statistically significant finding would be 1 - $0.95^{66} = 1 - 0.0339 = 0.966$.
In other words, students are likely to ``discover'' falsely positive results based on the large number of statistical tests involved.
We need to note, however, that there might well be real associations health and Zodiac sign, for health outcomes related to seasonality of birth (see for example \cite{carm:2022} and \cite{schiz:2005}.


Another activity that might spark discussion and reflection is working through the interactive visualization ``Hack Your Way to Scientific Glory" provided by \href{https://web.archive.org/web/20250117051842/http://projects.fivethirtyeight.com/p-hacking}{Five Thirty Eight}, where students play the role of a social scientist ``with a hunch" about the association between who is in office (Republicans vs.\ Democrats) and the state of the economy.
Students could be assigned into groups, each one taking on the exploration of this question from the vantage of a (randomly assigned) political party.
By choosing from a set of what politicians to include (e.g., Presidents, Governors, Senators), various measures of how to measure the economy (e.g., Employment, Stock Market), and choice of other factors (e.g., inclusion of recessions) students are able to see how analytic choices affect the p-value for the estimated slope coefficient in a regression model.
Students can share their ``best" model after settling on their choices.
This will often lead to a variety of models that yield statistically significant results ($p < 0.05$).
\cite{p-hacker} provides a similar environment to explore the perils of multiple testing.


\textbf{Assessment:} This activity can be assessed on its own in the form of replicability checklists and reflective writing and/or discussion about challenges encountered and implications for the progress of knowledge in their discipline.
Or, it can be assessed as part of selection and planning steps for a larger project of conducting a replication study.
A major component of the analysis plan for a replication study is a detailed review of the prior study's procedures, data, and availability thereof.

\textbf{Possible Extensions:}
One possible extension for students would be to use an applet or carry out a simulation study to repeatedly sample from a null distribution and to explore the behavior of confidence intervals for a desired population parameter (e.g., a difference in treatment group means).
See \href{https://www.rossmanchance.com/applets/2021/confsim/ConfSim.html}{here} or \href{https://sas-and-r.blogspot.com/2012/10/example-104-multiple-comparisons-and.html}{here} for examples without use of code and with code, respectively.

Another possible extension for students able to carry out simulations might involve assessing the Type I error-rate performance of a bivariate procedure to screen for variables to include in a multiple regression model.
Such a procedure (sometimes used in the literature) might then include those ``significant" predictors in a multivariate regression model.
Students could carry out a simulation study for m = 25 predictors when there were no significant predictors and observe the distribution of the overall test of the multiple regression model.
If there is really no signal (all of the predictors are independent of the outcome), we would hope that the overall error rate would be equivalent to our desired $\alpha = 0.05$ level.
The procedure described above will generally yield significant findings, despite the lack of ``real" differences.

\subsection{Reflective Assessment} \label{sec-assessment}

\textbf{Learning Goals and Outcomes:} Self-assessment and peer-assessment assignments can help make student work more reproducible by teaching them to become more reflective about aspects of reproducibility. 
Well-designed assessment-based assignments require students to synthesize aspects of recognizing reproducible and replicable research with aspects of appraising and assessing dubious research practices. 
Self-assessment and peer-assessment give students the chance to develop a greater capacity to identify appreciate reproducible and replicable research, identify problem areas, and to think critically about how to prioritize areas for improvement. 
Effective assessment assignments can help students partition reproducibility concepts into smaller cognitive units, which could increase their recall of core ideas and allow them to put more complex idea into practice. 
Self-assessment and peer-assessment assignments build on the content knowledge from our previous sections, but their learning goals are less based on knowledge acquisition and more based on developing their capacity for critical reflection. 

\textbf{Scaffolding/Skills:} Self-assessment and peer-assessment assignments can be incorporated at almost any level of instruction, but the benefits are greater when students have knowledge about reproducible and replicable research that they can use to support critical reflection. 
At the introductory level, a reflective assignment should include as much structure as possible.
Instead of asking 4-5 open-ended questions, an effective worksheet might ask students to state whether specific reproducibility or replicability practices can be found in a partner's writing, or in their own. 
This kind of reflective exercise is well suited for a no-code or low-code classroom. 
At more advanced levels, increasingly open-ended reflections can be deployed. 

The primary example discussed below is intended for the intermediate level and would work best in a classroom with direct instruction on coding or programming practices \citep{pruim:2023}. 
Scaffolding for this lesson plan can focus on debugging and refactoring strategies well-suited to novices \citep{McConnel:Code_Complete}, such as ``rubber ducking" (talking through code) and ``the binary chop'' (breaking a problem down to make it easier to solve). 
Depending on the level of the class, refactoring concepts such as functional decomposition or ``code smells" \citep{Bryan:2018}, might be more appropriate.

\textbf{Implementation:}
If a classroom includes direct instruction on coding practices, a low-stakes, in-class assignment can use the idea of reproducible code as a frame for peer-reflection and self-reflection. 
The professor can display or distribute block code that has been deliberately composed to raise errors if someone attempted to run it on their computer. 
Some these errors should relate to file and path names, package management, or other reproducibility and replicability concepts discussed above. 
Once the code snippet has been distributed to each student (via GitHub gist, downloaded file, etc.), students are prompted to: 

\begin{enumerate}
  \item read through the code chunk
  \item run the code, line by line, using an appropriate interpreter or software
  \item debug any easily solved coding errors 
  \item isolate or locate examples of specific ``code smells" or areas of concern that do not affect code execution 
\end{enumerate}

If the class is of a more introductory nature, students can be furnished with a list of known issues and be prompted to those issues within the block of sample code.
In a more intermediate or advanced context, students can be encouraged to ``bug hunt" on their own. 
Individual reflection on these items can be followed by a ``pair and share" activity, where students work with one or two classmates to compare notes and prepare for a larger group discussion. 
Subsequent whole-class discussion can focus on consensus takeaways, and the instructor can synthesize ideas and raise or emphasize anything not brought up or under-appreciated by the students. 

Written reflection assignments can build on more general pedagogical strategies, reinforcing R\&R as a set of principles, practices, and habits of mind, in contrast to merely providing a checklist of correct behaviors or end states. 
Introducing R\&R by way of code-based concepts may be especially appealing because of the ease with which an instructor can create flawed code. 
The code itself can be customized to match a use case that students have already seen, building on their familiarity and leaving more time for critical assessment. 

\textbf{Assessment:}
A low-stakes, in-class activity like the one discussed above can be assessed by structuring the activity so that all students end up participating, by giving oral feedback that praises incisive thought, attention to detail, language of meta-cognition, etc. 
Oral feedback does not need to focus on ineffective or incorrect ideas raised by students, but an instructor may wish to prompt students to go further with their thinking, rephrase for clarity, etc. 
In the context of a higher-stakes assignment with a reproducibility assessment worksheet component, or for a standalone reproducibility assessment assignment, a professor's feedback can model effective reflection and convey core knowledge about reproducibility. 
Since one of the overarching goals of reflection assignments is effective meta-cognition, written feedback should differentiate between evidence of effective meta-cognition and overall quality of writing. 
Ideally, types of meta-cognition specific to reproducibility---such as process over product, seeking core knowledge about reproducible research, and effective meta-cognitive regulation (i.e., how to monitor and control one's thinking)---would also be differentiated, largely by praising specific, positive examples of various reflective modes. 

\textbf{Possible Extensions:}
At the most advanced levels (e.g., a senior seminar or an advanced methods course), self-assessment and peer-assessment worksheets can begin with a checklists or a series of Likert-scale questions and work up to a set of more open-ended, reflective questions. 
Code and no-code considerations can be intermingled.
At least one question should prompt students to articulate priorities for revision or categorize the issues they discover as short-term or long-term concerns. 
Peer-assessment and self-assessment can also be sequenced with one another, such that students see examples of each others' reflections in ways that may positively influence their own.

\section{Conclusion}\label{sec-conc}

Reproducibility and replicability are not just technical publishing requirements.
They are foundational principles that shape how research is conducted, shared, and evaluated.
In order to create work that is actually reproducible, researchers must adhere to R\&R from the inception of a research project all the way through to its publication.

As educators, we must also teach R\&R as a habitual practice rather than a static criterion so that our students are able to engage in transparent and replicable research that stands up to scrutiny.
R\&R practices should be foundational to teaching research methods starting from the introductory level for two primary reasons:
First, habits that make research more difficult or impossible to reproduce begin early, so teaching reproducibility from the start can help discourage these bad habits while making R\&R practices more routine.
Second, teaching R\&R principles and practices helps cultivate the values of open and substantive engagement with the research community as one conducts, shares, and evaluates research.

As we write this paper in 2025, technological and social trends are making R\&R even more crucial.
On the technological side, the data and computational methods used in research are becoming more complex.
XX ADD a sentence or two that acknowledges challenges of AI and that using these tools (in professional and academic spheres) is often being done in ways that are not replicable, or keeping replicability in MIND. XX
Simultaneously, computers themselves are making computing for research appear more user-friendly but also more obscure through layers of graphical user interfaces, cloud storage and computing, and artificial intelligence. 
On the social side, science is emerging from a crisis of replicability with renewed emphasis on openness and transparency in research and publishing.
The public is profoundly divided on trust and funding for research at the same time that the world has urgent need for interdisciplinary convergence research to address major challenges and opportunities posed by artificial intelligence, climate change, pandemics, and more.
Taken together, these trends suggest a crucial and urgent need to transform how students learn research methods from the most introductory level in preparation for the data, computational tools, social expectations of research, and global problems that they will face.
We see tremendous value in integrating principles of reproducibility and replication directly into the curriculum across disciplines.
Not only will this approach pay dividends long term as we train researchers from the very onset with transparent research practices, but the principles and strategies highlighted in this article also have pedagogical value that can reinforce core course content and train students on skills central to their future as either producers or consumers of scientific findings.

We develop principles to achieve this R\&R transformation in teaching research methods and provide example pedagogical activities to put these principles to practice.
The principles outlined in this paper frame how research organization, planning, and reflection can help students internalize R\&R as both technical skills and ethical commitments to an ongoing practice.
Students can start learning R\&R practices and principles through activities that require just a single session or homework in addition to their existing syllabus.
These include reviewing published research, organizing a project reproducibly, rotating a project among peers, pre-analysis planning, appraising dubious research practices, and reflecting on their work and that of their peers.
The principles and practices articulated above are meant to lay a foundation for students to conduct and publish their own replicable research, and even to formally reproduce or replicate existing studies. 
They can help prepare students to independently repeat and extend the studies described in research papers and to participate in the research community by contributing their own reviews of current knowledge through reproduction and replication reports.

It has taken many years to start the R\&R transformation in science, including fundamental development of new research infrastructure and changes to incentive structures and the culture of science. 
It will likewise require broad shifts in educational infrastructure and culture of higher education to achieve an R\&R transformation in teaching research methods.
We urgently need the instructors of methods courses to start teaching R\&R one lesson at a time.
There is no need to reinvent the wheel: new instructors can stand on the shoulders of researchers and other educators who are publishing their own studies and educational materials with open licenses. 
In turn, we encourage instructors to publish original R\&R teaching materials as open educational resources to share.
The changes we have proposed here start at the grassroots level in the classroom, but more transformational change will also require integration of R\&R into departmental curricula and institutional learning goals.
Finally, the literature on pedagogical effectiveness of teaching R\&R remains largely anecdotal, suggesting a large need for reproducible and replicable education research on R\&R pedagogy.


\section{Acknowledgments}

\if1\anon{
The Alliance to Advance Liberal Arts Colleges funded a Reproducibility and Replication workshop which took place at Middlebury College in the summer of 2024 (see \url{10.17605/OSF.IO/GZHC8}, \cite{aalac_osf}). 
Holler and Kedron acknowledge support of National Science Foundation SBE 2049837.

Also contributing to this paper are participants in the workshop: [in alphabetical order]: Xavier Haro-Carrión (Geography, Macalester College), Emmanuel Kaparakis (Quantitative Analysis Center, Wesleyan University), Joseph J. Merry (Sociology, Furman University), and Anne Nurse (Sociology and Anthropology, The College of Wooster).
} \fi

\section{Disclosure statement}\label{disclosure-statement}

The authors do not have any conflicts of interest to declare.

\section{Data Availability Statement}\label{data-availability-statement}

We provide a compendium of lessons with advice for instructors at \url{https://doi.org/10.17605/OSF.IO/JQNHG}. 

\bibliography{bibliography.bib}

\end{document}
